\chapter{A Stiffness-Informed Robin Coupling}
\label{ch:method}

We replace the Dirichlet condition on the fluid with a Robin condition that combines velocity
and traction with a coefficient $\alpha$. Choosing $\alpha$ close to the wall's effective
impedance makes each sub-iteration nearly exact; we estimate it element by element from the
local wall thickness and stiffness, following the analysis of \citet{badia2008}.

\begin{figure}[htbp]
  \centering
  \begin{tikzpicture}[node distance=5mm and 8mm,
      box/.style={draw=accent, rounded corners=2pt, minimum height=9mm, align=center, font=\small}]
    \node[box] (f) {Fluid solve\\(Robin interface)};
    \node[box, right=of f] (s) {Structure\\solve};
    \node[box, right=of s] (c) {Converged?};
    \draw[-{Stealth}] (f) -- (s); \draw[-{Stealth}] (s) -- (c);
    \draw[-{Stealth}] (c.south) -- ++(0,-0.5) -| node[pos=0.25, below, font=\footnotesize] {no} (f.south);
  \end{tikzpicture}
  \caption{One time step of the partitioned Robin-Neumann scheme.}
  \label{fig:loop}
\end{figure}
